\documentclass[10pt,letterpaper]{article}
\usepackage{cornell-notes}

\setDocTitle{CISSP CBK Cornell Notes}
\setDocSubtitle{Domain 2: Asset Security}
\setDocVersion{Sixth Edition \textbullet\ 2021}
\setDocStatus{Study Companion}
\setDocOwner{Security Certification Study Collection}
\setDocAuthor{Arthur Deane and Aaron Kraus}
\setDocDate{\today}
\setDocSummary{Cornell-format study and review notes for Domain 2 of the CISSP CBK reference.}

\setCornellCollection{CISSP CBK Cornell Notes}
\setCornellUnitType{Domain}
\setCornellUnitNumber{2}
\setCornellUnitTitle{Asset Security}
\setCornellSourceLabel{CBK}
\setCornellSourceTitle{THE OFFICIAL (ISC)\textsuperscript{2} CISSP CBK REFERENCE}
\setCornellSourceAuthor{Arthur Deane and Aaron Kraus}
\setCornellSourceEdition{Sixth Edition \textbullet\ 2021}
\setCornellSourceRange{pp. 97--144}
\setCornellCompanionLabel{Study and review companion}

\begin{document}
\makecornelltitle

\section*{Domain Overview}
Asset security follows information and systems across their entire lifecycle: identify what exists, classify it by sensitivity/criticality/value, assign accountable roles, define handling rules, retain only as required, destroy defensibly, and apply controls appropriate to data state and organizational requirements.

\begin{summarybox}{Review Objectives}
Be able to distinguish classification from categorization; explain marking, handling, storage, and declassification; manage asset inventory and ownership; assign data lifecycle roles; plan retention and sanitization; distinguish data states; scope and tailor control baselines; select standards; and explain DRM/IRM, DLP, and CASB protections.
\end{summarybox}

\section{Identify and Classify Information and Assets}
\source{pp. 97--103}
\begin{cornell}
\noterow{Why inventory before protection?}{An organization cannot consistently protect assets it cannot identify or locate. Inventory and classification establish ownership, business value, sensitivity, criticality, location, and minimum control expectations. They also support risk assessment, incident response, compliance, and lifecycle decisions.}
\noterow{Data classification}{Organize data by sensitivity, criticality, or value so controls can be proportional to potential impact. The chapter describes \textbf{context-based} classification from metadata such as location/ownership, \textbf{content-based} classification from the information itself, and \textbf{user-based} classification assigned manually by knowledgeable users.}
\noterow{Commercial classification example}{A private-sector scheme may use levels such as confidential, private, sensitive, and public. Labels are organization-specific; the key is to define each level, ownership, handling rules, and protection expectations consistently.}
\noterow{Government classification cue}{Government/national-security schemes are more formal and jurisdiction-dependent. Do not assume an organization's commercial labels are equivalent to official government classifications or handling caveats.}
\noterow{Data classification vs. categorization}{\textbf{Classification} assigns sensitivity/criticality/value labels. \textbf{Categorization} groups information types with comparable sensitivity/impact so similar protections can be selected consistently. NIST SP 800-60 is cited for information-type categorization.}
\noterow{Asset classification}{Extends classification beyond data to systems, media, devices, services, and other assets that store, process, or transmit information. The classification of the information usually drives the minimum protection required for the containing asset.}
\noterow{Benefits of classification}{Enables proportional controls, prioritization, consistent handling, efficient spending, clearer accountability, improved incident decisions, and evidence of regulatory/contractual due care. Classification should be policy-driven and periodically reviewed.}
\end{cornell}

\section{Information and Asset Handling Requirements}
\source{pp. 104--107}
\begin{cornell}
\noterow{Marking and labeling}{Make classification and handling requirements visible to people and systems. Physical markings and digital labels may include classification, owner, distribution restrictions, and retention information. DLP and other automated tools can use digital labels.}
\noterow{Handling}{Define how each classification may be accessed, copied, transmitted, transported, shared, and used. Handling rules should apply to people, endpoints, removable media, networks, cloud services, and physical records as relevant.}
\noterow{Storage}{Select safeguards based on classification and risk. Sensitive digital data commonly requires access control and encryption at rest, while keys require protection at least commensurate with the data they protect. Physical media requires controlled storage and environmental protection.}
\noterow{Declassification}{Lower an asset's assigned sensitivity when authorized conditions are met. The data owner is central to the decision; document approvals and update markings, access, storage, and handling requirements. Failure to declassify can waste resources, while premature declassification creates exposure.}
\noterow{De-identification and tokenization}{De-identification reduces the ability to associate data with an individual. Tokenization substitutes a nonsensitive token for a sensitive value and keeps the mapping separately protected. Unlike ordinary encryption, a token is not mathematically reversed to recover the original; recovery uses the mapping/token vault.}
\end{cornell}

\section{Provision Resources Securely}
\source{pp. 108--114}
\begin{cornell}
\noterow{Information and asset ownership}{Every important information asset should have an accountable owner who determines acceptable use, classification, and protection requirements. Technical execution may be delegated, but accountability remains with the appropriate owner.}
\noterow{What belongs in an asset inventory?}{At minimum capture a unique identity, type, owner/custodian, location, business purpose, classification, responsible administrators, lifecycle/status, and dependencies. For technology assets include relevant hardware/software/configuration and data handled.}
\noterow{How can assets be discovered?}{Use procurement and configuration records, network/system discovery, software-license tools, DLP discovery, cloud inventories, endpoint management, and manual reconciliation. No single source is guaranteed complete; compare sources and address unmanaged assets.}
\noterow{Asset management lifecycle}{Provision through approved processes, establish secure configuration and ownership, track movement/change, maintain inventory and controls, review usefulness/authorization, and retire/dispose securely. Changes in ownership or location should update inventory and access.}
\noterow{Due care in asset management}{Reasonable protection includes accurate inventory, valid classification, defined safeguards, and appropriate handling. Due diligence verifies that these measures exist and continue to operate across mobile, cloud, and third-party environments.}
\end{cornell}

\section{Manage the Data Lifecycle}
\source{pp. 115--126}
\begin{cornell}
\noterow{Data owner}{Determines how and why data may be used, its classification, and required protection; authorizes access and remains accountable even when operational tasks are delegated. All sensitive data should have an identifiable owner.}
\noterow{Data controller}{Determines the purposes and means of processing personal data. In GDPR-oriented terminology, the controller must demonstrate compliance with processing principles and establishes contractual expectations for processors. A controller and owner can overlap but are not inherently identical.}
\noterow{Data custodian}{Technical role that maintains data and the supporting IT environment according to owner/business requirements. Custodians implement safe storage, transport, backups, permissions, and related operational safeguards.}
\noterow{Data processor}{Processes, transfers, or otherwise handles data on behalf of a controller/owner. A processor follows the controller's instructions and contractual/data-processing agreement; it does not decide the processing purpose in the same way the controller does.}
\noterow{Data user and data subject}{A \textbf{user} is an intended consumer/recipient of data. A \textbf{data subject} is the identifiable natural person about whom personal data is collected. Do not confuse the person using information with the person the information describes.}
\noterow{Collection / creation}{Apply privacy-by-design, secure defaults, classification, labeling, ownership, and access restrictions when data enters the environment. Decisions made here influence protection throughout the lifecycle. Collect only data that has a legitimate purpose and known requirements.}
\noterow{Data location}{Know where data is physically and logically stored/processed, including cloud regions, backups, replicas, endpoints, and third parties. Location can determine jurisdiction, residency obligations, contractual restrictions, and who controls encryption keys.}
\noterow{Data maintenance / use}{Keep classifications, permissions, quality, metadata, and safeguards current as data is accessed, updated, replicated, and shared. Balance usability with security; reassess controls as use and sensitivity change.}
\noterow{Data retention}{Keep information long enough to satisfy legal, regulatory, contractual, operational, evidentiary, and historical needs -- but not indefinitely without a reason. A retention schedule defines record type, owner, trigger/start date, retention period, and disposition.}
\noterow{Why destroy data?}{Expired/unneeded data still creates breach, discovery, privacy, storage, and maintenance exposure without business value. Destruction should follow an approved schedule, preserve legal holds, and produce appropriate evidence when certification is required.}
\noterow{What is data remanence?}{Residual information remains recoverable after inadequate deletion. Ordinary file deletion often removes references rather than every underlying data representation. Sanitization method must match media type, classification, regulatory duties, reuse plan, and required assurance.}
\noterow{Clear vs. purge vs. destroy}{\textbf{Clearing}: logical overwrite/wipe, the lowest assurance of the three. \textbf{Purging}: stronger sanitization such as degaussing where applicable. \textbf{Destruction}: render media unusable, e.g., shredding, pulverizing, or burning; the chapter also recognizes strong cryptographic destruction. Document the result.}
\noterow{Cloud destruction challenge}{Customers often cannot address the physical drive beneath an abstraction. Contracts, provider processes/evidence, lifecycle controls, tenant isolation, and cryptographic key control become central to meeting destruction/remanence requirements.}
\end{cornell}

\begin{exambox}{Key Distinction}
The owner/controller makes business or processing decisions; the custodian/processor performs operational handling under those requirements. Delegating work does not automatically delegate accountability.
\end{exambox}

\section{Ensure Appropriate Asset Retention}
\source{pp. 127--130}
\begin{cornell}
\noterow{What determines retention?}{Legal and regulatory mandates, contracts, business/operational value, audit needs, litigation/legal holds, historical value, privacy/minimization duties, and cost/risk. There is no universal retention period for all organizations or records.}
\noterow{Retention schedule}{Create a defensible schedule mapping records to owners, authoritative copies, origin/event dates, retention periods, and final disposition. Review it with legal/compliance and data/application owners on a recurring basis.}
\noterow{Retention best practices}{Label records; preserve accessible, authentic copies; address email and backups; deduplicate unnecessary copies; migrate from obsolete/degrading media; review schedules; delete expired data securely; and keep lifecycle handling controls in force until final disposition.}
\noterow{Legal hold vs. scheduled deletion}{A legitimate preservation requirement can suspend normal destruction. Ensure automated deletion honors legal holds and that release of a hold returns affected records to the approved retention/disposition process.}
\end{cornell}

\section{Data Security Controls and Compliance Requirements}
\source{pp. 131--144}
\begin{cornell}
\noterow{Technical, administrative, physical controls}{\textbf{Technical} controls use hardware/software capabilities; \textbf{administrative} controls shape human/process behavior through policy, standards, procedures, and governance; \textbf{physical} controls protect facilities and tangible assets. Layer categories to cover the risk.}
\noterow{Three data states}{\textbf{At rest}: stored and not actively processed/transmitted. \textbf{In transit/in motion}: moving between locations or systems. \textbf{In use}: actively processed, often in RAM/caches/registers. Controls must address each state.}
\noterow{Protecting data at rest}{Use access control plus appropriate encryption. Full-disk/volume, self-encrypting drive (SED), Trusted Platform Module (TPM)-supported protection, file-level, and field-level encryption offer different granularity and threat coverage.}
\noterow{Protecting data in transit}{Use trusted/secure protocols such as TLS and VPNs. \textbf{Link encryption} protects each communication link with decryption/re-encryption at intermediate nodes; \textbf{end-to-end encryption} protects content between communicating endpoints so intermediaries need not possess content keys.}
\noterow{Protecting data in use}{Strong identification/authentication, authorization, least privilege, application/process isolation, monitoring/accounting, and secure memory/system design are important. Storage encryption alone cannot stop a user/process that is legitimately decrypting and using the information.}
\noterow{Scoping vs. tailoring}{\textbf{Scoping} determines which baseline controls and assets are applicable. \textbf{Tailoring} modifies the selected baseline for organizational characteristics and risk. Removed/modified controls require risk-based justification and documentation -- convenience alone is not sufficient.}
\noterow{Common controls}{Controls whose protection is inherited by multiple systems/assets, such as enterprise physical entry, shared network monitoring, PKI, or organization-wide policy. Central implementation can reduce duplication if inherited effectiveness is actually verified.}
\noterow{Compensating controls}{Alternative/additional safeguards used when a prescribed control cannot feasibly meet its objective in the standard manner. The compensating control should provide equivalent risk mitigation, and both the reason and equivalence must be documented.}
\noterow{How should standards/frameworks be selected?}{Start with mandatory law/regulation/contract, then consider asset sensitivity, industry, jurisdiction, risk tolerance, business mission, and cost/benefit. The chapter surveys DoD RMF, NIST RMF/CSF, UK guidance, NIST 800-53/53A, FIPS 199/200, and ISO 27001/27002.}
\noterow{DRM vs. IRM}{Digital rights management (DRM) controls use, modification, copying, and distribution of digital intellectual property. Information rights management (IRM) applies similar persistent access/use restrictions more broadly to organizational files, email, and other information.}
\noterow{What is DLP?}{Data loss prevention discovers/classifies sensitive information, monitors its movement/use, and enforces policy to prevent accidental or intentional disclosure. Coverage can be storage/endpoint/network based; tune for false positives/negatives and account for encrypted traffic.}
\noterow{DLP lifecycle}{\textbf{Discovery/classification} identifies sensitive instances; \textbf{monitoring} observes relevant storage, endpoints, applications, and exfiltration paths; \textbf{enforcement} alerts, logs, quarantines, or blocks based on policy and classification.}
\noterow{DLP by data state}{At rest: discover/monitor repositories and complement encryption/access control. In transit: inspect network flows, subject to encrypted-traffic visibility limitations. In use: endpoint/host DLP can control actions such as removable-media transfer, copy/paste, or capture.}
\noterow{What is a CASB?}{A cloud access security broker sits logically between cloud consumers and services to centralize security policy enforcement. The chapter highlights four functions: \textbf{visibility}, \textbf{data security}, \textbf{threat protection}, and \textbf{compliance}.}
\end{cornell}

\section{Critical Comparisons}
\begin{center}\small
\begin{tabularx}{\textwidth}{>{\bfseries\color{Navy}\raggedright\arraybackslash}p{0.23\textwidth}>{\raggedright\arraybackslash}X>{\raggedright\arraybackslash}X}
\toprule Concept & First idea & Distinguishing idea\\\midrule
Classification / categorization & Assign sensitivity/criticality/value & Group information types with comparable impact/protection needs\\
Owner/controller / custodian/processor & Decide purpose, use, classification, protection & Implement custody or processing under those decisions\\
At rest / in transit / in use & Stored & Moving / actively processed\\
Scoping / tailoring & Decide applicability & Modify applicable control baseline for context\\
Clear / purge / destroy & Logical wipe & Stronger sanitization / render media unusable\\
DRM / DLP & Control permitted use of protected content & Detect and prevent sensitive-data leakage\\
\bottomrule
\end{tabularx}
\end{center}

\section{Key Terms and Acronyms}
\begin{summarybox}{Rapid-Review Glossary}
\textbf{CASB} -- Cloud Access Security Broker \quad
\textbf{DLP} -- Data Loss Prevention \quad
\textbf{DRM} -- Digital Rights Management \quad
\textbf{IRM} -- Information Rights Management \quad
\textbf{SED} -- Self-Encrypting Drive \quad
\textbf{TPM} -- Trusted Platform Module \quad
\textbf{Data owner} -- accountable decision maker for data use/protection \quad
\textbf{Custodian} -- technical keeper of data \quad
\textbf{Controller} -- determines processing purpose/means \quad
\textbf{Processor} -- processes for the controller \quad
\textbf{Remanence} -- recoverable residual data after inadequate sanitization.
\end{summarybox}

\section{High-Yield Review}
\begin{itemize}
\item Inventory and classification precede consistent control selection.
\item The most sensitive data on an asset often drives that asset's minimum protection.
\item Owners remain accountable when custodial/processing tasks are delegated.
\item Retention should be long enough to satisfy requirements but no longer than justified.
\item Deletion is not necessarily sanitization; select clear, purge, or destroy based on risk and media.
\item Protect data at rest, in transit, and in use with controls appropriate to the threat.
\item Scope determines applicability; tailoring adapts the baseline and requires documented rationale.
\item DLP = discover/classify, monitor, enforce. CASB = cloud visibility, data security, threat protection, compliance.
\end{itemize}

\section{Active-Recall Questions}
\begin{enumerate}
\item How does classification differ from categorization?
\item Who is ultimately accountable for defining data protection: owner or custodian?
\item Distinguish a data controller from a data processor.
\item Why should classification occur at or near collection/creation?
\item What is data remanence, and why is ordinary deletion insufficient?
\item Put clearing, purging, and destruction in increasing order of assurance.
\item What are the three data states and a primary protection concern for each?
\item How does scoping differ from tailoring?
\item When is a compensating control appropriate?
\item What are the three core DLP stages?
\item What four CASB functions does the chapter emphasize?
\item Why can excessive retention be a security problem?
\end{enumerate}

\subsection*{Answer Key}
\begin{enumerate}
\item Classification assigns sensitivity/criticality/value; categorization groups information types with similar impact/protection needs.
\item The owner; the custodian implements technical handling under the owner's requirements.
\item The controller determines purposes/means; the processor performs processing on the controller's behalf.
\item Early classification lets controls, labels, ownership, and access follow the data throughout its lifecycle.
\item It is recoverable residual data; ordinary deletion may remove pointers without eliminating the underlying representation.
\item Clearing, purging, destruction.
\item At rest/storage and access/encryption; in transit/secure transport; in use/authorization, isolation, and monitoring.
\item Scoping decides applicability; tailoring modifies the selected baseline for organizational context.
\item When the required control objective applies but the prescribed implementation is infeasible; document equivalent mitigation.
\item Discovery/classification, monitoring, enforcement.
\item Visibility, data security, threat protection, compliance.
\item It increases breach, privacy, discovery, storage, and maintenance exposure without corresponding business value.
\end{enumerate}

\section{Cornell Summary}
\begin{summarybox}{Domain 2 Summary}
Asset security converts knowledge about information into lifecycle protection. Inventory establishes what exists and who is accountable; classification and categorization establish value and sensitivity; handling and data-state controls preserve CIA; retention keeps records only as long as justified; and sanitization prevents obsolete assets from carrying recoverable data into reuse or disposal. Effective control baselines are scoped, tailored, documented, and supported by technologies such as rights management, DLP, and CASB where they fit the risk.
\end{summarybox}

\section{Final Domain Checklist}
{\small
\begin{itemize}[label=\(\square\),nosep]
\item I can distinguish data classification, categorization, and asset classification.
\item I can explain marking, handling, storage, declassification, and tokenization.
\item I can maintain ownership and inventory across the asset lifecycle.
\item I can distinguish owner, controller, custodian, processor, user, and subject.
\item I can build retention and secure-destruction requirements and explain remanence.
\item I can protect data at rest, in transit, and in use.
\item I can distinguish scoping, tailoring, common, supplemental, and compensating controls.
\item I can explain framework-selection criteria and DRM/IRM, DLP, and CASB roles.
\end{itemize}
}
\vfill
{\footnotesize\color{Teal}Prepared as a study companion to Deane \& Kraus, \textit{The Official (ISC)\textsuperscript{2} CISSP CBK Reference}, 6th ed. Content is paraphrased and synthesized for study.}

\end{document}
